\section{Monotone limits of analytic envelopes}
\label{sec:analytic-envelope-limits}

Demailly regularization preserves a curvature lower bound up to a vanishing error. Along dyadic indices, subadditivity also orders its analytic singularity types. The corresponding envelopes therefore decrease. Divisorial estimates identify the limit with the ideal envelope, by the same ideal-theoretic mechanism used in \cite[Theorem~7.1.1, proof of $(3)\Rightarrow(1)$, Step~1]{SGPT}.

Fix $u\in\PSH(X,\theta)$ and a Hermitian form $h$ on $X$. Choose the regularization from \cref{thm:demailly-hermitian-input} along $m_j=2^j$, starting at a sufficiently large index. Denote it by $a_j$ and choose numbers $\epsilon_j>0$ decreasing to zero such that
\begin{equation}\label{eq:dyadic-analytic-backgrounds}
 a_j\in\PSH(X,\theta+\epsilon_jh),\qquad
 a_j\geq u,\qquad
 a_j\text{ has type }(1/m_j,\mathcal I(m_ju)).
\end{equation}
The metric $h$ remains fixed, and the chosen regularization satisfies $a_j\to u$ in $L^1$. The values $a_j$ need not decrease. Their ordinary singularity types do: subadditivity in \eqref{eq:trace-local-ideal-inputs} gives
\[
 \mathcal I(2m_ju)\subseteq\mathcal I(m_ju)^2,
 \qquad a_{j+1}\preceq a_j.
\]
Taking logarithms of generators proves the ordinary comparison, with a constant allowed to depend on $j$.

\begin{theorem}\label{thm:dyadic-analytic-envelope-limit}
For every $u\in\PSH(X,\theta)$ and every choice in \eqref{eq:dyadic-analytic-backgrounds}, put
\[
 q_j=P_{\theta+\epsilon_jh}[a_j].
\]
Then $q_j$ has analytic singularities, and
\[
 q_j=P_{\theta+\epsilon_jh}[a_j]_{\mathcal I}
       \searrow P_\theta[u]_{\mathcal I}
\]
pointwise and in $L^1$. The limit is independent of the regularizing representatives, the admissible error sequence and the metric $h$.
\end{theorem}
\begin{proof}
By \cref{prop:analytic-envelope-type}, $q_j=P_{\theta+\epsilon_jh}[a_j]_{\mathcal I}\sim a_j$. Since the backgrounds decrease and $a_{j+1}\preceq a_j$, the function $q_{j+1}$ is a competitor for $P_{\theta+\epsilon_jh}[a_j]$. Hence $q_{j+1}\leq q_j$.

Put $p=P_\theta[u]_{\mathcal I}$. The inequality $a_j\geq u$ implies $u\preceq_{\mathcal I}a_j$. Regard $p$ in the larger background and use ideal projection and order preservation there. This gives
\[
 u-\sup_Xu\leq p\leq
 P_{\theta+\epsilon_jh}[p]_{\mathcal I}\leq q_j\leq0.
\]
Thus the decreasing limit $q$ is genuine, has supremum zero and belongs to $\PSH(X,\theta)$: the current inequalities pass to the limit, and a nested intersection of the compact maximum sets preserves the supremum. The displayed lower bound also proves $L^1$ convergence by dominated convergence.

\Needspace{6\baselineskip}
Fix a local divisor $E$. The inequalities $u-\sup_Xu\leq q\leq q_j$ imply
\[
 \nu_E(u)\geq\nu_E(q)\geq\nu_E(q_j)
       =\frac1{m_j}\operatorname{ord}_E\mathcal I(m_ju).
\]
By \cref{lem:bergman-divisorial-error}, the last term tends to $\nu_E(u)$. Therefore $q$ and $u$ have the same Lelong number on every local divisor over $X$. This gives the pointwise data on every local modification as well: blow up the point and use that its Lelong number is the generic number on the exceptional divisor; on a curve the point already is a divisor. The local ideal criterion \cref{thm:local-i-inputs} now gives $q\sim_{\mathcal I}u$.

Since $q\leq0$ is $\theta$-plurisubharmonic, it is a competitor for $p=P_\theta[u]_{\mathcal I}$. The earlier inequality $p\leq q$ proves equality. This identifies the limit solely from $\theta$ and the ideal class of $u$, proving every asserted independence.
\end{proof}

\Needspace{17\baselineskip}
\begin{corollary}\label{cor:dyadic-goodness-criterion}
For the sequence above, the following are equivalent:
\[
 P_\theta[u]=P_\theta[u]_{\mathcal I},\qquad
 q_j\longrightarrow P_\theta[u]\text{ in }L^1,
 \qquad q_j\searrow P_\theta[u]\text{ pointwise}.
\]
They hold for one admissible choice of approximation if and only if they hold for every such choice. For the fixed volume $\mathrm{d}V$,
\[
 \lim_j\int_X(q_j-P_\theta[u])\,\mathrm{d}V
   =\int_X(P_\theta[u]_{\mathcal I}-P_\theta[u])\,\mathrm{d}V\geq0,
\]
and equality to zero is another equivalent condition.
\end{corollary}
\begin{proof}
The theorem identifies the pointwise and $L^1$ limit. Equality almost everywhere is equality everywhere by \cref{lem:qpsh-almost-everywhere-order}. Finally $0\leq q_j-P_\theta[u]\leq-(u-\sup_Xu)$ almost everywhere, so dominated convergence gives the integral formula. A nonnegative difference has zero integral precisely when the two canonical representatives agree.
\end{proof}

\subsection{Absorbing the curvature error}

For a Hermitian background $\omega$, take $\theta=h=\omega$ in \eqref{eq:dyadic-analytic-backgrounds}. A rational rescaling absorbs the curvature error and preserves the rational coefficients of analytic singularities. It gives an approximation entirely within $\PSH(X,\omega)$.

\Needspace{14\baselineskip}
\begin{corollary}\label{cor:fixed-hermitian-analytic-density}
Let $\omega$ be Hermitian and let $u\in\PSH(X,\omega)$. There are analytic $b_j\in\PSH(X,\omega)$ with $b_j\to u$ in $L^1$, with decreasing ordinary singularity types, such that
\[
 P_\omega[b_j]=P_\omega[b_j]_{\mathcal I}\searrow P_\omega[u]_{\mathcal I}.
\]
Consequently analytic singularity classes are dense in the entire ideal-type metric space of \cref{prop:ideal-lone-metric}. The equality $P_\omega[u]=P_\omega[u]_{\mathcal I}$ is equivalent to convergence of these ordinary envelopes to $P_\omega[u]$ in $L^1$.
\end{corollary}
\begin{proof}
Use $h=\omega$ in \eqref{eq:dyadic-analytic-backgrounds}. Enlarge the error bounds to positive rational numbers $\epsilon_j\downarrow0$, and set
\[
 c_j=\frac1{1+\epsilon_j},\qquad b_j=c_ja_j.
\]
Then $c_j\uparrow1$ is rational and
\[
 \omega+\ddc b_j
 =c_j\bigl((1+\epsilon_j)\omega+\ddc a_j\bigr)\geq0.
\]
The analytic coefficient is $c_j/m_j$. Also $b_j\to u$ in $L^1$, since $a_j\to u$ there and the sequence has bounded $L^1$ norm. Because $a_{j+1}\preceq a_j$, $c_{j+1}\geq c_j$, and each $a_j$ is bounded above, one has $b_{j+1}\preceq b_j$.

Put $q_j=P_{(1+\epsilon_j)\omega}[a_j]=P_{(1+\epsilon_j)\omega}[a_j]_{\mathcal I}$. The preceding theorem gives $q_j\searrow P_\omega[u]_{\mathcal I}$ pointwise and in $L^1$. Since $b_j$ is analytic, \cref{prop:analytic-envelope-type} and the homogeneity in \cref{prop:i-envelope-operations} give
\[
 r_j:=P_\omega[b_j]=P_\omega[b_j]_{\mathcal I}=c_jq_j.
\]
Here $q_j\leq0$ decreases and $c_j\uparrow1$, so $r_j$ also decreases, with pointwise limit $P_\omega[u]_{\mathcal I}$. Moreover
\[
 P_\omega[u]_{\mathcal I}\leq q_j\leq r_j\leq0.
\]
The genuine lower bound proves $L^1$ convergence by dominated convergence.

The metric distance from $[b_j]$ to $[u]$ is exactly the $L^1$ distance from $r_j$ to $P_\omega[u]_{\mathcal I}$, so it tends to zero. The last assertion follows from the identified limit.
\end{proof}
